% 第2章 经济学地基
\section{缝隙存在的经济学地基}\label{sec:economics}

本章不引入新结论，只做一件铺垫：说清为什么“缝隙”这种东西应该存在、又为什么它在逻辑上必然有限。三段地基——有效市场的悖论结构、信息进入价格的机制、低信噪比的压迫——决定了后文每条命题的形式与语气。

\subsection{有效市场之争真正留下的结论}

Fama（1970）\cite{fama1970} 把有效市场分成弱式、半强式、强式三档，把“价格是否已反映某类信息”变成了可检验命题。半个世纪的争论留下的结论不是“市场有效”，而是一个更微妙的东西：\textbf{价格对新信息的反应快到让大多数慢半拍的人赚不到钱；但快的人能赚，而快是有成本的。}这句话已经埋下了本文的全部素材——缝隙存在，但它的规模被成本与速度锁在一个有限的区间里。

Grossman 与 Stiglitz（1980）\cite{grossmanstiglitz1980} 把这个微妙之处做成了严格的逻辑。若价格已充分反映所有信息，就没人有动力花成本获取信息；可一旦没人获取，价格又不可能充分反映信息。市场只能停在一个不动点上：价格足够有效，使获取信息的边际成本恰等于边际收益——\textbf{而这个均衡里必然留着一部分未被消除的偏差}。推论很强：任何“市场已经（或将被）完全有效”的说法都不成立，均衡总把不可消除的残余留给成本最低的获取者。

残余为什么不被立刻套掉？Shleifer 与 Vishny（1997）\cite{shleifer1997} 给出了机制：套利有资金约束，也承担噪声交易者风险——偏差可以继续扩大，套利者可能等不到均值回归就被迫平仓。所以偏差能长期存在，不是无人看见，而是看见的人受资金与耐心约束。这段文献的含义是：\textbf{缝隙的存在与大小，从提出之日起就是一个关于“谁能以更低成本看见并承受”的问题，而不是关于“价格对不对”的问题。}

\subsection{信息进入价格的机制：租金自动耗散}

Kyle（1985）\cite{kyle1985} 给出了信息进入价格的最简机制。市场里有一类知情交易者掌握私有信息，一类噪声交易者随机下单，做市商只能看到净订单流，从中间接推断信息。均衡下，价格以速率
\begin{equation}
  \lambda = \frac{1}{2}\sqrt{\frac{\Sigma_0}{\sigma_u^2}}
\end{equation}
向真实价值收敛，其中 $\Sigma_0$ 是私有信息的先验方差，$\sigma_u^2$ 是噪声交易的方差。这个式子说的是：\textbf{知情交易者越多、信息越有价值，价格对订单流越敏感，单位信息能换取的利润就越薄。}这是一种自动的租金耗散——与 Grossman--Stiglitz 讲的是同一件事，只是从订单流的层面重讲了一遍。

对本文而言，Kyle 的 $\lambda$ 有一个具体角色：它是第 \ref{sec:measurement} 章“信息衰减速率”的微观基础之一。缝隙的尺度窗口向高频端收缩还是向低频端展开，机制上由这类价格冲击系数与信息扩散速度共同决定。另有一个经验侧面：McLean 与 Pontiff（2016）\cite{mcleanpontiff2016} 发现，因子在被论文发表后收益显著衰减，且把发表后交易活动的变化考虑进来后衰减更深。租金耗散不是理论假想，它在数据里看得见。

\subsection{低信噪比对方法论的压迫}

Gu、Kelly 与 Xiu（2020）\cite{gukellyxiu2020} 用一整套机器学习方法预测美股横截面收益，报告的样本外 $R^2$ 只有约 $0.3\%$--$0.4\%$（月度）。这个数字单看小得可笑，换算成多空组合的年化收益与夏普却可以相当可观。它揭示的道理是：\textbf{金融里“解释力低”与“能赚钱”并不矛盾——杠杆与横截面聚合能放大微弱信号；但反过来，“样本内 $R^2$ 高”也几乎不能说明任何问题，任何一个过参数化模型都能把噪声拟合得很漂亮。}

低信噪比带来三个连锁后果。其一，过拟合的风险被放到最大，样本内每一次调参都是一次数据窥探\cite{white2000}，任何候选池都自带向上偏差。其二，模型比较的统计功效极低，两个模型的差异常常落在噪声里，用“谁的回测曲线更好看”来选模型，等于在噪声里选。其三，也是本文最常回到的一条：\textbf{信号的真实性必须以“经济上可交易”来检验，而不是以统计指标来检验。}第 \ref{sec:definition} 章的定义、第 \ref{sec:criteria} 章的统计门槛，都是这条压迫的产物。

\subsection{两个相对性的预告}

后文两个核心命题——尺度窗口与参照系分解——在地基里各有源头，这里先各交代一句，免得它们出现在后文时显得突兀。

尺度一侧：Clark（1973）\cite{clark1973} 提出价格的从属过程（subordinated process）模型，价格在“信息时间”而非日历时间上展开，交易活跃度就是这两种时间之间的变换；Müller 等（1993）\cite{muller1993} 把异质性推到更强：市场不是一个人的市场，而是多个时间尺度上参与者的叠加，不同尺度上观测到的“规律”本就可以呈现不同形态。更早，Working（1960）\cite{working1960} 证明了一个技术性但根本的事实——对随机过程做时间聚合，会改变其可观测的自相关结构。三点合起来指向同一句话：\textbf{“在日线上看到的”与“在月线上看到的”，不是同一现象的两个样本，而是两个不同的问题。}

参照系一侧：Roll（1977）\cite{roll1977} 指出市场组合不可观测，CAPM 的一切检验都悬在“你用的是什么代理基准”上；Fama（1970）\cite{fama1970} 的联合假设问题说得更直接——任何有效性检验都是“市场有效”与“定价模型正确”的联合检验，无法单独证伪其一。翻译成本文的语言：\textbf{$\alpha$ 不是收益自身的属性，而是收益与参照系的联合属性。}

\subsection{本章要点}

缝隙的存在有经济学解释：它是信息租金在成本与约束下未被耗尽的残余，其数学形态受制于租金耗散机制。残余有多大，一半取决于信息本身，另一半取决于观测它的装置——这正是接下来两条线各自要推进的内容：第 \ref{sec:measurement} 章管时间尺度，第 \ref{sec:reference} 章管参照系。